\section{The model problem}

We consider a reaction-diffusion problem governed by the conservation law
\begin{equation}
	\nabla \cdot \sP + \zeta\,\uP = q
	\label{eq:conservation}
\end{equation}
posed on a regular domain $\Omega \subset \mathbb{R}^{d_{\Omega}}$. In eq.~(\ref{eq:conservation}) the symbol $\nabla$ represents the vector of partial derivatives, so that $\nabla \cdot \sP$ is the divergence of the vector field $\sP$, which represents the flux. Also, $\uP$ is the concentration of the conserved quantity, $\zeta \in L^{\infty}(\Omega)$ is the (possibly position-dependent) reaction coefficient and $q\in L^2(\Omega)$ is the source. In thermal applications $\uP$ is the temperature, $\sP$ is the heat flux and $q-\zeta \uP$ is a heat source that depends linearly on the temperature. In weakly compressible flow in porous media $\uP$ is the pressure and $\sP$ is the fluid flux.

The boundary of $\Omega$, denoted by $\partial \Omega$, is divided into a Dirichlet part $\partial \Omega_D$, where $\uP=\uP_D$ is imposed, and a Neumann part $\partial\Omega_N$, in which $\sP \cdot \mathbf{\check{n}}=H$ is imposed. Hereafter, we assume for simplicity homogeneous boundary conditions throughout, i.e., $\uP_D = 0$ and $H = 0$.

Let us also define the field $\eP$ as the gradient of $\uP$,
\begin{equation}
	\eP = \nabla \uP
\label{eq:compatibility}
\end{equation}

This equation can be viewed as a compatibility condition. We call a vector field $\eP$ compatible if there exists a scalar field $\uP$ such that $\eP=\nabla \uP$.

Now, suppose that two vector fields $\stP$ and $\etP$ are given in $\Omega$. These fields will both be assumed in $\left[L^2(\Omega)\right]^{d_{\Omega}}$ and thus possibly discontinuous. Since they originate from a measured dataset, it is generally the case that $\stP$ does not satisfy the conservation law, and $\etP$ does not satisfy the compatibility condition. Moreover, they are concomitant, in the sense that each sample $\etP$ corresponds to the same spatial location as $\stP$ -- that is, they form coupled observations of gradient and flux at the same point in the domain --.

The goal, therefore, is to determine the best approximations to $\stP$ and $\etP$ within the space of feasible fields, that is, those that are both conservative and compatible. The scalar field $\uP$ emerges naturally as a by-product of this reconstruction.

\subsection{Well-posedness of the Exact Problem}

We begin by introducing the function spaces involved in the formulation:

\begin{equation}
	\left.
	\begin{array}{rcl}
		U & = & \left\{ \bar{u} \in H^1(\Omega),\ \bar{u} = 0 \text{ on } \partial\Omega_D \right\} \\
		E & = & \left[ L^2(\Omega) \right]^{d_\Omega} \\
		S & = & \left[ L^2(\Omega) \right]^{d_\Omega} \\
	\end{array}
	\right \rbrace  \quad X := U \times E \times S
\end{equation}
We now pose the problem as a minimization over the feasible set $ Z \subset X $, consisting of all tuples $ (\uP, \eP, \sP) \in X $ that satisfy the governing equations and boundary conditions.
\begin{tcolorbox}[colback=black!10!white, halign = justify, boxrule=0pt]
\begin{problem*}[P1]
	Find $x= \left (\uP,\eP,\sP \right) \in Z$ -- the feasible space -- that minimizes the distance to the data in the feasible space, i.e.,
	\begin{equation}
		J(x) \leq J(x+\delta x)~,\qquad \forall \,\delta x\,\in\,\delta Z~
	\end{equation}
	where 
	\begin{equation}
		J(\bar{\uP},\bar{\eP},\bar{s})=\frac12 \|\bar{\eP}-\etP\|^2 + \frac{\kappa}{2} \|\bar{\sP}-\stP\|^2~,
	\end{equation}
	and the feasible space is
	\begin{equation}
		Z=\left \{ (\bar{\uP},\bar{\eP},\bar{\sP})\in X~|~\nabla \cdot \bar{\sP} + \zeta \,\bar{\uP} = q~,~\bar{\eP}=\nabla \bar{\uP},\,\left.~\bar{\sP}\cdot\mathbf{\check{n}}\right |_{\partial \Omega_N}=0 \right \}
	\end{equation}
	with the variation given by
	\begin{equation}
		\delta Z=\left \{ (\delta{\uP},\delta{\eP},\delta{\sP})\in X~|~\nabla \cdot \delta{\sP} + \zeta \delta{\uP} = 0~,~ \delta{\eP}=\nabla \delta{\uP} ,\,\left. ~\delta{\sP}\cdot\mathbf{\check{n}}\right |_{\partial \Omega_N}=0\right \}~.
	\end{equation}
\end{problem*}
\end{tcolorbox}

\begin{remark}
	Above, $\|\cdot\|$ is the norm of $\left[L^2(\Omega)\right]^{d_{\Omega}}$. The corresponding scalar product is denoted by $\left ( \cdot , \cdot \right )$. The same notation will be used for $L^2(\Omega)$.
\end{remark}
Let us now relax the constraints through Lagrange multipliers $\lP$ and $\mP$, which belong to the following function spaces:
\begin{equation*}
	L=\{g\in H^1(\Omega)~|~g=0 \mbox{ on } \partial \Omega_D\},\qquad
	\mbox{ and }\qquad M=\left[L^2(\Omega)\right]^{d_{\Omega}},  
\end{equation*}
respectively. We also use the notation
\begin{equation}
	y=\left (\lP,\mP\right )\,\in \, Y:=\,L\times M~.
\end{equation}
Then, the resulting equivalent problem can be written as: 
\begin{tcolorbox}[colback=black!10!white, halign = justify, boxrule=0pt]
	\begin{problem*}[P2]
		Find $(x,y)=\left( \uP, \eP, \sP, \lP, \mP \right) \in X\times Y$ such that 
		\begin{equation}
			(x,y)= \mbox{arg} \left \{ \inf_{\bar{x}\in X} \sup_{\bar{y}\in Y} 
			\mathcal{L}(\bar{x},\bar{y})
			\right \}
		\end{equation}
		where
		\begin{equation}
			\begin{aligned}
				\mathcal{L}(\bar{x},\bar{y}) &~=~J(\bar{\uP},\bar{\eP},\bar{\sP})+ \nonumber \\
				&~-~\left ( \bar{\sP},\nabla \bar{\lP} \right ) + \left (\zeta \bar{\uP} - q , \bar{\lP} \right )  
				+ \left ( \bar{\eP}-\nabla \bar{\uP}, \bar{\mP} \right ) ~.
			\end{aligned} \label{eq:lagrangian}
		\end{equation}
	\end{problem*}
\end{tcolorbox}
The weak form of the optimality conditions reads:
\begin{equation}\label{eq:variational_form}
	\begin{array}{rcll}
		\left ( \lP, \zeta \delta \uP \right ) - \left ( \mP, \nabla \delta \uP \right ) & = & 0 &\qquad \forall \,\delta \uP \in U \\
		\left ( \eP, \delta \eP \right ) + \left ( \mP, \delta \eP \right ) &=& \left ( \etP, \delta \eP \right )  &\qquad \forall \,\delta \eP \in E\\
		\kappa \left ( \sP, \delta \sP \right ) - \left ( \nabla \lP, \delta \sP \right )  & = & \kappa \left ( \stP,\delta \sP \right ) &\qquad \forall\, \delta \sP \in S \\
		\left ( \zeta \,\uP, \delta \lP \right )-\left ( \sP, \nabla \delta \lP \right )  &=& \left ( q, \delta \lP \right ) 	&\qquad \forall\,\delta \lP \in L\\
		-\left (\nabla \uP,\delta \mP \right ) + \left ( \eP, \delta \mP \right ) & = & 0 &\qquad \forall \, \delta \mP \in M 
	\end{array}
\end{equation}

The sum of the left-hand sides (with the last two equations multiplied by $-1$) defines the bilinear form of the problem in the product space $X \times Y$, given by:

%\begin{eqnarray}
% \left ( \lP, \zeta\delta u \right ) - \left ( \mP, \nabla \delta u \right ) & = & 0 \qquad\qquad\qquad\qquad \forall \,\delta u\in U  \label{eq:var1}\\
%\left ( e,\delta e \right )+ \left ( \mP, \delta e \right ) &=& \left ( \etP,\delta e \right )  \qquad\qquad\qquad \forall \,\delta e \in E\label{eq:var2}\\
%\kappa \left ( s, \delta s \right ) - \left ( \nabla \lP, \delta s \right )  & = & \kappa \left ( \stP,\delta s \right ) \qquad\qquad\qquad \forall\, \delta s \in S \label{eq:var3}\\
%\left ( \zeta\,u,\delta \lP \right )-\left ( s,\nabla \delta \lP\right )  &=&\left ( q,\delta \lP \right ) \qquad\qquad\qquad \forall\,\delta \lP \in L\label{eq:var4}\\
%-\left (\nabla u,\delta \mP \right )+\left ( e,\delta \mP \right ) & = & 0 \qquad\qquad\qquad\qquad \forall\,\delta \mP \in M \label{eq:var5}
%\end{eqnarray}
\begin{equation}
	B((x,y),(\delta x,\delta y))=a(x,\delta x)+b(\delta x,y)-b(x,\delta y)
\end{equation}
where
\begin{equation}\label{eq:bform}
	\begin{array}{rcl}
		a(x,\delta x)&=& \left (\eP,\delta \eP \right )+\kappa \left ( \sP,\delta \sP \right ) \\
		b(x,y)& = & \left ( \mP,\eP-\nabla \uP \right )-\left (\nabla \lP,\sP\right ) + \left ( \lP, \zeta \uP \right )
	\end{array}
\end{equation}

The optimality conditions in strong form, therefore, are given by:
\begin{eqnarray}
	\zeta \lP  + \nabla \cdot \mP & = & 0 			\label{eq:dif1}\\
	 \eP + \mP & = & \etP 							\label{eq:dif2}\\
	\kappa\, \sP - \nabla \lP & = & \kappa\,\stP 	\label{eq:dif3}\\
	 \zeta\,\uP+\nabla \cdot \sP  & = & q 			\label{eq:dif4}\\
	-\nabla \uP + \eP &=& 0 						\label{eq:dif5}
\end{eqnarray}
together with
\begin{eqnarray}
 \sP \cdot \mathbf{\check{n}} & = & 0~,\qquad \mbox{ on } \partial \Omega_N \\
 \mP \cdot \mathbf{\check{n}} & = & 0~,\qquad \mbox{ on } \partial \Omega_N ~.
\end{eqnarray}

\begin{remark}
	Notice that when the reaction coefficient $\zeta$ is identically zero the five variational equations given in (\ref{eq:variational_form}) -- or, equivalently, their differential versions, i.e., equations (\ref{eq:dif1}) to (\ref{eq:dif5}) -- are in fact two independent sets of equations: (a) The first, second and fifth equations, with unknowns $\uP$, $\eP$ and $\mP$, and (b) the third and fourth equations, with unknowns $\sP$ and $\lP$. In such a case, $\sP$ is the best approximation to $\stP$ among conservative flux fields, and $\eP=\nabla \uP$ is the best approximation to $\etP$ among compatible (gradient) fields, and both problems can be analyzed separately. The reaction term adds an interesting twist to the problem, coupling all five variational equations.	
\end{remark}

\begin{theorem}[Well-posedness of exact problem]
\label{theorem1}
	Problem P2 is well posed.
	
	\begin{proof}
		Following the theory of mixed problems \cite{brezzi-fortin,ern-guermond} one needs to show that $a(\cdot,\cdot)$ and $b(\cdot,\cdot)$ are continuous, that $a(\cdot,\cdot)$ is coercive in $\delta Z$, with
		\begin{eqnarray}
			\delta Z&=&\{x\in X~|~b(x,y)=0\, \quad \forall y\in Y\}\nonumber \\
			&=&\{(\uP,\eP,\sP)\in X \mid \eP-\nabla \uP=0, ~\nabla \cdot \sP + \zeta \uP =0,~\left. ~\sP \cdot \mathbf{\check{n}}\right |_{\partial\Omega_N} = 0  \}~,
		\end{eqnarray}
		and that the inf-sup condition on $b(\cdot,\cdot)$,
		\begin{equation}
			\inf_{y\in Y}\sup_{x\in X} \frac{b(x,y)}{\|x\|_X\|y\|_Y} \geq k_b > 0~.
		\label{eqinfsup}
		\end{equation}
		holds true.
		
		The norms adopted in $X$ and $Y$ are: 
		\begin{eqnarray}
			\|x\|_X &=& \|(\uP,\eP,\sP)\|_X = \sqrt{\|\nabla \uP\|^2+\|\eP\|^2+\kappa \|\sP\|^2}~,\\
			\|y\|_Y &=& \|(\lP,\mP)\|_Y=\sqrt{\frac{1}{\kappa}\left \| \nabla \lP \right \|^2+\left\|\mP \right \|^2}~.
		\end{eqnarray}
		
		They were chosen so that they have the same units. It is easy to see that both $X$ and $Y$ are Hilbert spaces, and that $a(\cdot,\cdot)$ and $b(\cdot,\cdot)$ are continuous. Notice that
		\begin{equation*}
				a(x,x)=\|\eP\|^2 + \kappa \|\sP\|^2 ~.
		\end{equation*}
		So, for $x \in \delta Z$ it holds that:
		\begin{equation}
			a(x,x)= \frac12 \|\eP\|^2 +\frac12 \|\nabla \uP\|^2+ \kappa \|\sP\|^2 \geq \frac12 \|x\|_X^2~.
		\label{eqcoercivity}
		\end{equation}
		and, thus, coercivity holds with constant equal to $\frac12$.
		
		It remains to show that there exists $k_b>0$ such that:
		\begin{equation*}
			\sup_{(\uP,\eP,\sP)\in X} \frac{\left (\mP,\eP-\nabla \uP \right ) - \left (\nabla \lP, \sP \right ) + \left ( \lP, \zeta \uP \right )}{\|(\uP,\eP,\sP)\|_X} \geq k_b \, \|(\lP,\mP)\|_{Y}~, \qquad \forall (\lP,\mP)\in Y~.
		\end{equation*}
	
		To this end, we can take $\uP=0$, $\eP=\mP$ and $\sP=-\frac{1}{\kappa}\nabla \lP$ to get
		\begin{equation*}
			\sup_{(\uP,\eP,\sP)\in X} \frac{\left (\mP,\eP-\nabla \uP \right ) - \left (\nabla \lP, \sP \right ) + \left ( \lP, \zeta\,\uP \right )}{\|(\uP,\eP,\sP)\|_X} 
			\geq \frac{\|\mP\|^2 + \frac{1}{\kappa}\|\nabla \lP\|^2}{\sqrt{\|\mP\|^2 + \frac{1}{\kappa}\|\nabla \lP \|^2}} = \|(\lP,\mP)\|_Y~,
		\end{equation*}
		giving an inf-sup constant $k_b$ equal to one. Thus, the coercivity of $a(\cdot,\cdot)$ on $\delta Z$ and the inf-sup stability of $b(\cdot,\cdot)$ have been proven. Then we have the well-posedness of the exact problem, with a unique solution $(\uP,\eP,\sP)\in H^1(\Omega) \times \left[L^2(\Omega)\right]^{d_{\Omega}} \times L^2(\Omega)$ that depends continuously on $q\in L^2(\Omega)$, $\etP \in \left[L^2(\Omega)\right]^{d_{\Omega}}$ and $\stP \in \left[L^2(\Omega)\right]^{d_{\Omega}}$.
	\end{proof}
\end{theorem}


\subsection{Approximation}

Let $U_h$, $E_h$, $S_h$, $L_h$ and $M_h$ be the finite element spaces for the five unknowns. We assume that $U_h\subset U$ and $L_h\subset L$ (in particular, the Dirichlet boundary conditions are strongly enforced). 
Also, we assume that $E_h$, $S_h$ and $M_h$ are subspaces of $[L^2(\Omega)]^d_{\Omega}$ and define $X_h:=U_h\times E_h \times S_h$ and $Y_h:=L_h\times M_h$. The numerical approximation we propose corresponds to solving the following discrete problem:

\begin{tcolorbox}[colback=black!10!white, boxrule=0pt, halign=justify]
	\begin{problem*}[P2$_{h}$]
		Find $(x_h,y_h)=\left ( \uP_h, \eP_h, \sP_h, \lP_h, \mP_h\right ) \in X_h\times Y_h$ such that:
		
		\begin{equation}
			(x_h,y_h)= \mbox{arg} \left \{ \inf_{\bar{x}\in X_h} \sup_{\bar{y}\in Y_h} 
			\tilde{\mathcal{L}}(\bar{x},\bar{y})
			\right \}
		\end{equation}
		where
		\begin{eqnarray}
			\tilde{\mathcal{L}}(\bar{x},\bar{y})&=&J(\bar{\uP},\bar{\eP},\bar{\sP}) \nonumber \\
			&-&\left ( \bar{\sP},\nabla \bar{\lP} \right ) + \left (\zeta \bar{\uP} - q , \bar{\lP} \right )  
			+ \left ( \bar{\eP}-\nabla \bar{\uP}, \bar{\mP} \right ) \nonumber \\
			&+&\frac{\alpha}{2} \left \| \bar{\eP} - \nabla \bar{\uP} \right \|^2 \nonumber \\
			&-&\frac{\gamma}{2} \|\bar{\eP}+\bar{\mP}-\etP\|^2 \nonumber \\
			&-&\frac{\eta}{2\kappa} \|\kappa \bar{\sP} - \nabla \bar{\lP} - \kappa \stP\|^2 \nonumber \\
			&+& \frac{\theta \kappa \ell_{\sP}^2 }{2}\, \|\nabla \cdot \bar{\sP} + \zeta\, \bar{\uP} - q \|_h^2 \nonumber \\
			&-& \frac{\beta \ell_{\mP}^2 }{2}\, \|\nabla \cdot \bar{\mP} + \zeta\, \bar{\lP} \|_h^2~.
		\end{eqnarray}
	\end{problem*}
\end{tcolorbox}
\noindent
in which we kept the saddle-point structure by adding the squared norms of the optimality conditions for the dual variables and subtracting the square norms of the optimality conditions for the primal ones.

%\noindent{\bf Discrete Problem P2$_h$:} Find $(x_h,y_h)=\left ( u_h, e_h, s_h, \lP_h, \mP_h\right ) \subset X_h\times Y_h$ such that 
%\begin{equation}
%(x_h,y_h)= \mbox{arg} \left \{ \inf_{\bar{x}\in X_h} \sup_{\bar{y}\in Y_h} 
%\tilde{\mathcal{L}}(\bar{x},\bar{y})
%\right \}
%\end{equation}
%where
%\begin{eqnarray}
%\tilde{\mathcal{L}}(\bar{x},\bar{y})&=&J(\bar{u},\bar{e},\bar{s})+ \nonumber \\
%& &-\left ( \bar{s},\nabla \bar{\lP} \right ) + \left (\zeta \bar{u} - q , \bar{\lP} \right )  
%+ \left ( \bar{e}-\nabla \bar{u}, \bar{\mP} \right ) +\nonumber \\
%& &+\frac{\alpha}{2} \left \| \bar{e} - \nabla \bar{u} \right \|^2 + \\
%& &-\frac{\gamma}{2} \|\bar{e}+\bar{\mP}-\etP\|^2 + \nonumber \\
%& &-\frac{\eta}{2\kappa} \|\kappa \bar{s} - \nabla \bar{\lP} - \kappa \stP\|^2 \nonumber \\
%& &+ \frac{\theta \kappa \ell_{\sP}^2 }{2}\, \|\nabla \cdot \bar{s} + \zeta\, \bar{u} - q \|_h^2 + \nonumber \\
%& & - \frac{\beta \ell_{\mP}^2 }{2}\, \|\nabla \cdot \bar{\mP} + \zeta\, \bar{\lP} \|_h^2~.
%\end{eqnarray}

The mesh-dependent norm $\|\cdot\|_h$ is defined as 
%\begin{equation}
$\|f\|_h^2 = \sum_K \|f\|^2_{L^2(K)} $
%\end{equation}
with the sum performed over all elements $K$ of the mesh. The associated scalar product is denoted as $\left ( \cdot, \cdot \right )_h$. This mesh-dependent norm is only necessary when the finite element spaces for $\mathbf{s}$ or $\boldsymbol{\mu}$ are discontinuous, since their divergence is then not in $L^2(\Omega)$. For continuous approximations of $\mathbf{s}$ and $\boldsymbol{\mu}$ the norm $\|\cdot\|_h$ coincides with the $L^2(\Omega)$-norm.

Above, $\alpha$, $\gamma$, $\eta$, $\theta$ and $\beta$ are non-negative non-dimensional coefficients to be further discussed shortly, and $\ell_{\sP}$ and $\ell_{\mP}$ are length scales. The function $\tilde{\mathcal{L}}$ is an augmented Lagrangian, of which the saddle point is the solution of P2$_h$. 

The optimality conditions of P2$_h$ are:

\begin{eqnarray}
	\alpha \left ( \nabla \uP_h, \nabla \delta \uP \right ) + \theta \kappa \ell_{\sP}^2  \left ( \zeta^2 \uP_h, \delta \uP \right ) - \alpha \left ( \eP_h,\nabla \delta \uP \right ) + & & \nonumber \\ 
	 + \theta \kappa \ell_f^2  \left ( \nabla \cdot \sP_h, \zeta \delta \uP \right )_h + \left ( \lP_h, \zeta\delta \uP \right ) - \left ( \mP_h, \nabla \delta \uP \right ) & = & \theta \kappa \ell_{\sP}^2  \left ( q, \zeta\delta \uP \right )\nonumber \\ & & \quad \forall \,\delta \uP \in U_h  \label{eq:var1h} \\ 
	-\alpha \left ( \nabla \uP_h, \delta \eP \right ) +(1+\alpha-\gamma)\,\left ( \eP_h,\delta \eP \right )+ (1-\gamma)\, \left ( \mP_h, \delta \eP \right ) &=& \left (1-\gamma)\,( \etP, \delta \eP \right )  \nonumber \\ & &\quad \forall \,\delta \eP \in E_h\label{eq:var2h}\\
	(1-\eta)\kappa \left ( \sP_h, \delta \sP \right ) - (1-\eta)\left ( \nabla \lP_h, \delta \sP \right ) + & & \nonumber \\
	+\theta \kappa \ell_{\sP}^2  \left ( \nabla \cdot \sP_h, \nabla \cdot \delta \sP \right )_h+ \theta \kappa \ell_{\sP}^2 \,\left (\zeta\,\uP_h, \nabla \cdot \delta \sP \right )_h & = & (1-\eta)\kappa \left ( \stP,\delta \sP \right ) + \nonumber \\ & &+\theta \kappa \ell_{\sP}^2  \left ( q , \nabla \cdot \delta \sP \right )_h \nonumber \\ & & \quad \forall\, \delta \sP \in S_h \label{eq:var3h}\\
	-(1-\eta)\left ( \sP_h,\nabla \delta \lP\right ) + \left ( \zeta\,\uP_h,\delta \lP \right ) - \beta \ell_{\mP}^2 \left ( \zeta^2 \lP_h, \delta \lP\right )+ & & \nonumber \\
	- \beta \ell_{\mP}^2 \left ( \nabla \cdot \mP_h, \zeta\,\delta \lP \right )_h -\frac{\eta}{\kappa} \left ( \nabla \lP_h , \nabla \delta \lP \right )&=&\left ( q,\delta \lP \right ) + \eta \left ( \stP,\nabla \delta \lP \right ) \nonumber \\ & & \quad \forall\,\delta \lP \in L_h\label{eq:var4h} \\
	-\left (\nabla \uP_h,\delta \mP \right )+(1-\gamma)\left ( \eP_h,\delta \mP \right )
	- \beta \ell_{\mP}^2 \left ( \zeta\,\lP_h, \nabla \cdot \delta \mP \right )_h + & & \nonumber \\- \beta \ell_{\mP}^2 \left ( \nabla \cdot \mP_h, \nabla \cdot \delta \mP \right )_h - \gamma \left ( \mP_h,\delta \mP \right ) & = & -\gamma \left (\etP,\delta \mP \right ) \nonumber \\ & & \quad \forall\,\delta \mP \in M_h \label{eq:var5h}
\end{eqnarray}

The sum of the left-hand sides (the last two multiplied by $-1$) is nothing but the discrete bilinear form that corresponds to P2$_h$, i.e.,
\begin{equation}
\tilde{B}((x_h,y_h),(\delta x,\delta y))=\tilde{a}(x_h,\delta x)+\tilde{b}(\delta x,y_h)-\tilde{b}(x_h,\delta y)+\tilde{c}(y_h,\delta y)~,
\end{equation}
where
\begin{equation}
	\begin{array}{rcl}
		\tilde{a}(x_h,\delta x)~&=&~(1-\gamma)\left (\eP_h,\delta \eP\right )+(1-\eta)\kappa \left ( \sP_h,\delta \sP \right ) \\
		~&+&~\alpha \left( \eP_h-\nabla \uP_h, \delta \eP - \nabla \delta \uP \right)\\
		~&+&~ \theta \kappa \ell_{\sP}^2  \left ( \nabla \cdot \sP_h + \zeta \uP_h , \nabla \cdot \delta \sP + \zeta \, \delta \uP \right )_h~,
	\end{array}
\end{equation}

\begin{equation}
	\begin{aligned}
		\tilde{b}(\delta x,y_h)~=~b(\delta x,y_h)-\gamma \left (\mP_h,\delta \eP\right )+\eta \left (\nabla \lP_h,\delta \sP \right ) ~,
	\end{aligned}
\end{equation}

\begin{equation}
	\begin{array}{rcl}
		\tilde{c}(y_h,\delta y) ~&=&~  \beta \ell_{\mP}^2 \left ( \nabla \cdot \mP_h + \zeta \lP_h , \nabla \cdot \delta \mP + \zeta\,\delta \lP \right )_h \\
		~&+&~\dfrac{\eta}{\kappa}\,\left ( \nabla \lP_h,\nabla \delta \lP \right ) + \gamma \left ( \mP_h,\delta \mP \right )~,
	\end{array}
\end{equation}
and $b(\cdot,\cdot)$ as defined in (\ref{eq:bform}).

Problem P2$_h$ can thus be written as:

\begin{tcolorbox}[colback=black!10!white, boxrule=0pt, halign=justify]
	Find $(x_h,y_h)\in X_h \times Y_h$ such that:
	\begin{equation}
		\tilde{B}\left ( \left ( x_h,y_h \right ),\left ( \delta x,\delta y \right ) \right )= \tilde{F}(\delta x,\delta y), \quad \forall (\delta x, \delta y)\in X_h\times Y_h
		\label{eq:BzF}
	\end{equation}
\end{tcolorbox}

The linear form $\tilde{F}:X_h\times Y_h \to \mathbb{R}$ stands for the sum of the right-hand sides of (\ref{eq:var1h})-(\ref{eq:var5h}), with the last two multiplied by $-1$.

We now define the formulations considered in the analysis:

\begin{description}
\item[(a)] {\bf Natural formulation}. It corresponds to a choice of discrete spaces that satisfies
$$
\nabla U_h \subseteq M_h, \qquad \nabla L_h \subseteq S_h, \qquad E_h = M_h 
$$
When such is the case, as we will see, no stabilization is required. The coefficients are thus taken equal to zero,
$$
\alpha = \gamma = \eta = \theta = \beta = 0~.
$$
The typical spaces for this formulation is continuous piecewise $P_k$ for $U_h$ and $L_h$, and discontinuous, piecewise $P_{k-1}$ for $M_h$, $S_h$ and $E_h$.

\item[(b)] {\bf Equal-order formulation}. It corresponds to choosing the same $H^1$-conforming finite element space for all unknown fields, typically a continuous piecewise $P_k$ approximation. Equal-order formulations require stabilization. We consider three choices of the stabilization parameters:

\begin{description}

\item[(b-1)] {\bf Unstabilized method:}
All coefficients set to zero. This method is not optimally convergent.

\item[(b-2)] {\bf Minimally-stabilized method:} Only coefficients $\alpha$ and $\eta$ are positive. In particular, we take
$$
\alpha = \frac{1}{8}, \qquad \eta = \frac12, \qquad \gamma = \theta = \beta = 0~.
$$

\item[(b-3)] {\bf Fully-stabilized method:} All coefficients are positive, and the scale lengths are equal to the mesh size; i.e.,
	\begin{equation}
	\alpha = \frac{1}{8}, 
	\quad \gamma=\alpha, \quad \eta = \frac12,\quad \theta=\frac12,\quad \beta=\frac12,
	\quad \ell_{\sP} = \ell_{\mP} = h~. \label{eq:paramsopt}
	\end{equation}

\end{description}

\end{description} 

\noindent With these definitions, it is possible to prove

\bigskip

\begin{theorem}[Well-posedness of discrete problem and optimal approximation] For methods (a), (b-2) and (b-3) above, the discrete problem P2$_h$ is well-posed, and $(x_h,y_h)$ satisfies that $\exists c>0$ independent of $h$ such that
	\begin{equation}
		\||(x,y)-(x_h,y_h)\|| \leq c\,\min_{(\bar{x},\bar{y})\in X_h\times Y_h}
		\||(x,y)-(\bar{x},\bar{y})\||
		\label{eq:bestapprox}
	\end{equation}
where the norm $\||\cdot\||$ depends on the method: For methods (a) and (b-2) it is
\begin{equation}
\||(\bar{x},\bar{y})\||^2 ~=~ \|(\bar{x},\bar{y})\|^2_{X\times Y} ~=~ \|\nabla \bar{\uP} \|^2 +\|\bar{\eP}\|^2 + \kappa \|\bar{\sP}\|^2 +  \frac{1}{\kappa} \|\nabla \bar{\lP}\|^2 + \|\bar{\mP}\|^2 ~,
\end{equation}
and for method (b-3) it is 
\begin{equation}
\||(\bar{x},\bar{y})\||^2 ~=~ \|(\bar{x},\bar{y})\|^2_{X\times Y}  + \kappa h^2 \|\nabla \cdot \bar{\sP}\|^2  + h^2\,\|\nabla \cdot \bar{\mP}\|^2~.
	\end{equation}
\end{theorem}

%The extra terms in the Lagrangian add coercivity to the discrete problem and provide stability in the norm $\||\cdot\||_h$, allow us to work in a norm with improved control over $\nabla \uP_h$, $\nabla \cdot \sP_h$ and $\nabla \cdot \mP_h$.

\begin{proof}
	It is easy to verify that the exact solution $(x,y)$ satisfies the discrete problem (\ref{eq:BzF}) when substituted in the place of $(x_h,y_h)$. The formulation is thus consistent. The continuity of all the involved bilinear and linear forms in the norm $\||\cdot\||$ is straightforward. The rest of the proof is different for each method.
	
	\bigskip
	
 \noindent{\bf (a) Natural formulation}. It is the classical mixed approximation, in which the discrete bilinear forms are $\tilde{a}=a$ and $\tilde{b}=b$ and $\tilde{c}=0$. Let
  $$
 \delta Z_h=\{ x_h\in X_h ~|~b(x_h,y_h)=0\quad \forall y_h\in Y_h \}
 $$
be the discrete kernel of $b(\cdot,\cdot)$. Since $\nabla U_h\subset E_h=M_h$ we have $\eP_h = \nabla \uP_h$ and thus $\delta Z_h \subset \delta Z$. We thus have coercivity in the discrete kernel ($a(x_h,x_h) \geq \frac12 \|x_h\|_X^2$ for all $x_h \in \delta Z_h$) immediately from (\ref{eqcoercivity}). 

According to Brezzi's theorem \cite{Demkowicz2006} it only remains to show that the inf-sup condition (\ref{eqinfsup}) holds for the discrete spaces $X_h$ and $Y_h$ with a constant independent of the mesh. Consider an arbitrary $y_h = \left ( \lP_h, \mP_h \right )\in Y_h$ and take $\uP_h=0$, $\eP_h=\mP_h$ and $\sP_h = -\frac{1}{\kappa} \nabla \lP_h$. The two last choices are possible, since $E_h=M_h$ and $\nabla L_h \subset S_h$. One then gets
 \begin{equation}
 \sup_{x_h=(\uP_h,\eP_h,\sP_h)\in X_h} \frac{b\left (x_h,y_h\right )}{\|x_h\|_X} \geq \|y_h\|_Y~.
 \label{eqsupb}
 \end{equation}
 with inf-sup constant equal to 1.
% Strong coercivity in $\delta Z_h$, together with (\ref{eqsupb}), imply the claimed result (see Brezzi \& Fortin, 1991, Prop. 2.7, or Girault \& Raviart, 1986, Theorem 1.2).
 
 \bigskip
 
 \noindent{\bf (b-2) Equal-order minimally-stabilized method.} For this case we use Babuska's theorem \cite{Demkowicz2006}, which requires us to prove that $\exists K_B>0$, independent of $h$, such that
 \begin{equation}
 \sup_{(\bar{x},\bar{y})\in X_h\times Y_h} \frac{\tilde{B}\left ( (x_h,y_h), (\bar{x},\bar{y}) \right )}
 {\||(\bar{x},\bar{y})\||} \geq K_B\,\||(x_h,y_h)\||~,
 \end{equation}
 for all $(x_h,y_h)$ in $X_h\times Y_h$. Let $k_0 > 0$ and
 $(\bar{x},\bar{y})=\left (\uP_h, \eP_h - k_0 \mP_h, \sP_h, \lP_h, \mP_h \right )$. Then,
 \begin{eqnarray}
 \tilde{B}\left ((x_h,y_h), (\bar{x},\bar{y}) \right ) & = &
 \alpha \|\nabla \uP_h\|^2 + (1+\alpha) \|\eP_h\|^2 + (1-\eta) \kappa \|\sP_h\|^2 + \nonumber \\
 & & + k_0 \|\mP_h\|^2 + k_0 \alpha \left ( \nabla \uP_h, \mP_h \right ) 
 - k_0 (1+\alpha) \left ( \eP_h,\mP_h \right ) + \nonumber \\
 & & - 2\alpha \left ( \eP_h , \nabla \uP_h \right ) +\frac{\eta}{\kappa} \|\nabla \lP_h\|^2
 + \eta \left ( \eP_h , \mP_h \right )~.
 \end{eqnarray}
 Now we use Young's inequality ($2ab\geq -k a^2-(1/k)b^2$, $\forall a, b \in \mathbb{R}, \forall k>0$) as follows
 \begin{eqnarray*}
 -2\alpha \left ( \eP_h, \nabla \uP_h \right ) & \geq &
 -\frac{1}{k_1} \| \eP_h \|^2 - k_1 \alpha^2 \|\nabla \uP_h \|^2 \\
 \left ( \nabla \uP_h, \mP_h \right ) & \geq & -\frac{1}{2k_2} 
 - \frac{k_2}{2} \|\mP_h\|^2 \\
 -\left ( \eP_h, \mP_h \right ) & \geq & -\frac{1}{2 k_3} \|\eP_h \|^2
 - \frac{k_3}{2} \|\mP_h\|^2 \\
 \eta \left ( \eP_h, \mP_h \right ) & \geq & -\frac{\eta}{2 k_4} \|\mP_h\|^2
 - \frac{\eta k_4}{2} \|\eP_h\|^2
 \end{eqnarray*}
 to arrive at
 \begin{eqnarray*}
 \tilde{B}\left ((x_h,y_h), (\bar{x},\bar{y}) \right ) & \geq & 
 A_1 \|\nabla \uP_h\|^2 + A_2 \|\eP_h\|^2 + (1-\eta) \kappa \|\sP_h\|^2 +
 \frac{\eta}{\kappa} \|\lP_h\|^2 + A_3 \|\mP_h\|^2~,
 \end{eqnarray*}
 with
 \begin{eqnarray*}
 A_1 & = & \alpha - k_1 \alpha^2 - \frac{\alpha k_0}{2 k_2}~, \\
 A_2 & = & 1+\alpha - \frac{1}{k_1} - \frac{(1+\alpha)k_0}{2 k_3} - \frac{\eta k_4}{2}~, \\
 A_3 & = & k_0 - \frac{\alpha k_0 k_2}{2} - \frac{(1+\alpha)k_0k_3}{2} - \frac{\eta}{2 k_4}~.
 \end{eqnarray*}
 The above expressions are valid for any $k_i>0$, $i=0,\ldots,4$. Numerical optimization was used to determine the following set of parameters
 $$
 \alpha=\frac18,\quad \eta=\frac12,\quad k_0=0.352, \quad k_1=2.606, \quad
 k_2=0.606,\quad k_3=0.707, \quad k_4=1.654
 $$
 which produce strictly positive values $A_1 \simeq A_2 \simeq A_3 \simeq 0.047$.  At the same time, it is easy to see that, with $k_0=0.352$, $\||(\bar{x},\bar{y})\||\leq \sqrt{2}\,\||(x_h,y_h)\||$, which leads to
 $$
 K_B \geq \frac{0.047}{\sqrt{2}} \simeq 0.033~,
 $$
 independent of $h$ as required.
 
 \bigskip
 
 \noindent{\bf (b-3) Equal-order fully-stabilized method.} Similarly to (b-2), this case is proved using Babuska's theorem \cite{Demkowicz2006}. There is a difference, however, since when all stabilizing terms are active the bilinear form $\tilde{B}(\cdot,\cdot)$ is {\em strongly} coercive in the norm $\||\cdot\||$, i.e., $\exists K_B > 0$ (independent of $h$) such that
	\begin{equation}
		\tilde{B}((x_h,y_h),(x_h,y_h))\geq K_B \left ( 
		\|x_h\|_X^2 + \|y_h\|_Y^2 + \kappa\,h^2 \|\nabla \cdot \sP_h \|^2
		+ h^2\,\|\nabla \cdot \mP_h \|^2\right ),
	\end{equation}
	for all $ (x_h,y_h)\in X_h\times Y_h$~.
	By direct computation, 
	\begin{eqnarray}
		\tilde{B}((x_h,y_h),(x_h,y_h)) ~&=&~ (1-\gamma) \|\eP_h\|^2 + (1-\eta)\kappa \|\sP_h\|^2 \nonumber \\
		~&+&~ \alpha \|\eP_h-\nabla \uP_h\|^2 + \theta \kappa \ell_{\sP}^2 \|\nabla \cdot \sP_h + \zeta \uP_h\|^2 \nonumber \\
		~&+&~ \beta \ell_{\mP}^2 \| \nabla \cdot \mP_h + \zeta \lP_h \|^2
		+ \frac{\eta}{\kappa} \|\nabla \lP_h\|^2 + \gamma \|\mP_h\|^2~.
	\end{eqnarray}
	
	Now, as before, we apply Young's inequality to each of the cross-product terms in the previous expression.
	
	\begin{eqnarray}
		-2\alpha \left (\eP_h,\nabla \uP_h \right ) & \geq & -\frac14 \|\eP_h\|^2 - 4\alpha^2 \|\nabla \uP_h \|^2 \\
		%%%%%
		2\theta\kappa \ell_{\sP}^2\left ( \nabla \cdot \sP_h, \zeta \uP_h \right )_h& \geq & - \frac12 \kappa \theta \ell_{\sP}^2 \|\nabla\cdot \sP_h\|^2 - {2\theta\kappa\ell_{\sP}^2}\|\zeta \uP_h\|^2\\
		%%%%%%
		2 \beta \ell_{\mP}^2\left ( \nabla \cdot \mP_h, \zeta \lP_h \right )_h& \geq & - \frac12 {\beta} \ell_{\mP}^2 \|\nabla\cdot \mP_h\|^2 - {2\beta\ell_{\mP}^2}\|\zeta \lP_h\|^2
		%%%%%%%
	\end{eqnarray}
	
	Then, adding up all terms that contain $\sP_h$ we arrive at
	\begin{equation}
		(1-\eta) \kappa \|\sP_h\|^2 + \frac12 \kappa \ell_{\sP}^2 
		\theta  \|\nabla\cdot \sP_h\|_h^2~.
		\label{eq:term1}
	\end{equation}
	Adding up all terms that contain $\zeta \uP_h$ gives
	$$
	-\theta \kappa \ell_{\sP}^2 \|\zeta \uP_h\|^2
	$$
	which using Poincar\'e-Friedrichs inequality ($\|\uP_h\| \leq c_{\mbox{\scriptsize{PF}}}\,\ell_{\Omega}\,\|\nabla \uP_h\|$, 
	with $\ell_{\Omega}$ the diameter of $\Omega$) is shown to be greater or equal than 
	\begin{equation}
		- {\theta \kappa \ell_{\sP}^2 \|\zeta\|_{\infty}^2 c_{\mbox{\scriptsize{PF}}}^2 \ell_{\Omega}^2}\,\|\nabla \uP_h \|^2~.
		\label{eq:term2}
	\end{equation}
	By $\|\zeta\|_{\infty}$ we denote the $L^{\infty}(\Omega)$-norm of $\zeta$. Adding up all terms that contain $\mP_h$ we arrive at
	\begin{equation}
		\gamma \|\mP_h\|^2 +\frac12 \beta\ell_{\mP}^2  \|\nabla \cdot \mP_h \|_h^2~.
		\label{eq:term4}
	\end{equation}
	The sum of the terms that contain $\zeta \lP_h$ is greater or equal, by Poincaré-Friedrichs inequality, than
	\begin{equation}
		- \beta \ell_{\mP}^2 \|\zeta\|_{\infty}^2 c_{\mbox{\scriptsize{PF}}}^2 \ell_{\Omega}^2\,\|\nabla \lP_h \|^2~.
		\label{eq:term5}
	\end{equation}
	Adding up all terms that contain $\nabla \uP_h$ gives
	\begin{equation}
		\left ( \alpha - 4 \alpha^2 \right )
		\|\nabla \uP_h\|^2~.
		\label{eq:term6}
	\end{equation}
	Adding up all terms that contain $\eP_h$ gives
	\begin{equation}
		\left ( \frac34-\gamma + \alpha \right ) \|\eP_h\|^2~.
		\label{eq:term7}
	\end{equation}
	Adding everything up we get
	\begin{eqnarray}
		\tilde{B}((x_h,y_h),(x_h,y_h))~&\geq&~ \left (\frac34 -\gamma+\alpha \right ) \|\eP_h\|^2 + (1-\eta)\kappa \|\sP_h\|^2 + \nonumber \\
		~&+&~ (\alpha - 4 \alpha^2 - \theta \kappa \ell_{\sP}^2 \|\zeta\|_{\infty}^2 c_{\mbox{\scriptsize{PF}}}^2 \ell_{\Omega}^2)\, \|\nabla \uP_h\|^2 + \nonumber \\
		~&+&~\frac12 \theta \kappa \ell_{\sP}^2 \|\nabla \cdot \sP_h\|^2 + \nonumber \\
		~&+&~\frac12 \beta \ell_{\mP}^2 \| \nabla \cdot \mP_h\|^2 + \nonumber \\
		~&+&~\left ( \frac{\eta}{\kappa}-\beta \ell_{\mP}^2 \|\zeta\|_{\infty}^2 c_{\mbox{\scriptsize{PF}}}^2 \ell_{\Omega}^2 \right ) \|\nabla \lP_h\|^2 + \gamma \|\mP_h\|^2~. \label{eq:Bhparams}
	\end{eqnarray}
	Taking
	\begin{equation}
	\alpha = \frac{1}{8}, \quad \theta=\frac12,\quad \beta=\frac12,
	\quad \gamma=\alpha, \quad \eta = \frac12,
	\quad \ell_{\sP} = \ell_{\mP} = h, \label{eq:paramsopt}
	\end{equation}
	we get
	\begin{eqnarray*}
		\tilde{B}\left ( (x_h,y_h),(x_h,y_h) \right ) &\geq  &
		\left ( \frac{1}{16} - O(h^2) \right )\, \|\nabla \uP_h\|^2+\\
		& & +\frac{3}{4}\,\|\eP_h\|^2+\frac12 \kappa \|\sP_h\|^2 + \frac14 \kappa h^2 \|\nabla \cdot \sP_h\|_h^2 + \\
		& &
		+\left (\frac{1}{2\kappa}  -O(h^2) \right ) \|\nabla \lP_h\|^2 + \\
		& & + \frac18 \|\mP_h\|^2  +\frac14 h^2 \|\nabla \cdot \mP_h\|_h^2 
	\end{eqnarray*}
	and finally
	\begin{equation}
		\tilde{B}\left ( (x_h,y_h),({x}_h,{y}_h) \right ) \geq 
		K_B \, \| |(x_h,y_h)  \| |_h^2 
	\end{equation}
	where
	\begin{equation}
		K_B \geq \frac{1}{16}~.
	\end{equation}
	We have assumed that $h$ is sufficiently small so that the terms $O(h^2)$ can be neglected and thus $K_B$ becomes mesh independent. 
	
	\bigskip
	
	In all three cases we have proved that $\tilde{B}$ is consistent, continuous and coercive, which implies that the discrete problem is well posed and that the best approximation property (\ref{eq:bestapprox}) holds true.
\end{proof}


